\section{Entropy cost of a conditional centroid}\label{sec:entropy}
For $\tau>0$, let $\gamma_\tau$ be the centered Gaussian density on
$\R^d$ with covariance $\tau^2 I$. If $X$ is bounded, write
$f_X=\Law(X)*\gamma_\tau$ and
$\Ent(f)=-\int_{\R^d}f\log f$. All entropies and divergences in this
section use natural logarithms. Gaussian mixtures with bounded centers
have finite entropy and the Gaussian tail bounds needed below.

\begin{lemma}[Quadratic cost of barycentric compression]
\label{lem:centroid-entropy}
Let $X$ be a bounded random vector and $J$ a finite random variable on
the same probability space. Set $M=\E[X\mid J]$. Then
\begin{equation}\label{eq:centroid-entropy}
 \Ent(f_X)-\Ent(f_M)
       \le\frac{\E\norm{X-M}^2}{2\tau^2}.
\end{equation}
No sign assertion for the left side is required.
\end{lemma}
\begin{proof}
Put $f=f_M$. Direct differentiation of the finite Gaussian mixture gives
\begin{equation}\label{eq:log-hessian}
 \nabla^2\log f(y)
 =-\tau^{-2}I+\tau^{-4}\operatorname{Cov}(M\mid M+\tau G=y)
 \succeq-\tau^{-2}I,
\end{equation}
where $G$ is an independent standard Gaussian. This identity also follows
by differentiating the log-sum of exponentials after factoring out
$\exp(-\norm y^2/(2\tau^2))$.
Define $\varphi(x)=-\E_G\log f(x+\tau G)$. Boundedness of the centers
justifies differentiation and gives
$\nabla^2\varphi\preceq\tau^{-2}I$. Therefore
\[
 \varphi(X)-\varphi(M)
 \le \langle\nabla\varphi(M),X-M\rangle
           +\frac{\norm{X-M}^2}{2\tau^2}.
\]
The first term averages to zero conditional on $J$. On the other hand,
nonnegativity of $\KL(f_X\Vert f)$ gives
\[
 \Ent(f_X)-\Ent(f)
 \le -\int(f_X-f)\log f
 =\E\varphi(X)-\E\varphi(M).
\]
These inequalities prove the result. The Gaussian density used for
convolution is a common mathematical kernel; no additional noisy state
is assumed in the transducer.
\end{proof}

The Gaussian score/Hessian mechanism is classical. For comparison,
Polyanskiy--Wu \cite[Proposition 1 and Corollary 4]{PW} control smoothed
entropy differences by a Wasserstein distance. Here the specified
martingale coupling $M=\E[X\mid J]$ cancels the linear Taylor term.
The resulting one-sided \emph{squared} error, rather than a square-root
continuity estimate for unrelated measures, is what can be summed along
a stochastic register. We use the explicit proof above, not an assumed
entropy monotonicity under convex order.

\subsection{An entropy increment from a spherical gap}
Let $\Pi$ be spherical averaging at each radius on $L^2(\R^d)$:
\[
 (\Pi u)(r\omega)=\int_{\mathbb S^{d-1}}u(r\theta)\,d\sigma(\theta).
\]
It is the orthogonal projection onto radial functions. The angular gap
\eqref{eq:gap-intro}, integrated in polar coordinates, gives
\begin{equation}\label{eq:radial-gap}
 \norm{Pu}_2^2\le\norm{\Pi u}_2^2
                  +\lambda^2\norm{u-\Pi u}_2^2.
\end{equation}
Here $P$ acts on Euclidean functions by the same orthogonal pullbacks;
it preserves radial functions and commutes with $\Pi$.

\begin{lemma}[Entropy increment and angular Hellinger variance]
\label{lem:entropy-gap}
For a Gaussian mixture $f$ with bounded centers,
\begin{equation}\label{eq:entropy-gap}
 \Ent(Pf)-\Ent(f)
 \ge (1-\lambda^2)\norm{\sqrt f-\Pi\sqrt f}_2^2.
\end{equation}
Without a strict gap the entropy increment is still nonnegative.
\end{lemma}
\begin{proof}
Set $f_a(y)=f(U_a^{-1}y)$ and $g=\sum_ap(a)f_a=Pf$.
All $f_a$ have the same entropy, and hence
\[
 \Ent(g)-\Ent(f)=\sum_ap(a)\KL(f_a\Vert g).
\]
For probability densities $u,v$, Jensen's inequality gives
$\KL(u\Vert v)\ge-2\log\int\sqrt{uv}
\ge2(1-\int\sqrt{uv})=\norm{\sqrt u-\sqrt v}_2^2$.
Consequently the entropy increment is at least
\begin{align*}
 \sum_ap(a)\norm{\sqrt{f_a}-\sqrt g}_2^2
 &=1-\norm{P\sqrt f}_2^2
       +\norm{P\sqrt f-\sqrt g}_2^2\\
 &\ge1-\norm{P\sqrt f}_2^2.
\end{align*}
As $\norm{\sqrt f}_2^2=1$, \eqref{eq:radial-gap} proves the assertion.
The last statement also follows directly from entropy concavity. This
argument uses the full $L^2$ gap only on square roots of smoothed
probability densities; it assumes no logarithmic Sobolev inequality and
no one-step mixing of point masses.
\end{proof}
